\documentclass[12pt,letterpaper]{article}

% APA 7th Edition packages
\usepackage[utf8]{inputenc}
\usepackage[spanish]{babel}
\usepackage{amsmath}
\usepackage{amsfonts}
\usepackage{amssymb}
\usepackage{graphicx}
\usepackage{float}
\usepackage{array}
\usepackage{booktabs}
\usepackage{multirow}
\usepackage{multicol}
\usepackage{geometry}
\usepackage{setspace}
\usepackage{titling}
\usepackage{fancyhdr}
\usepackage{apacite}
\usepackage[colorlinks=true,linkcolor=black,citecolor=black,urlcolor=black]{hyperref}

% APA formatting
\geometry{margin=1in}
\doublespacing
\setlength{\parindent}{0.5in}

% Header configuration
\pagestyle{fancy}
\fancyhf{}
\rhead{\thepage}
\lhead{INFORME TÉCNICO TDC}
\renewcommand{\headrulewidth}{0pt}

% Title page information
\title{Análisis Térmico y de Transferencia de Calor:\\ Informe Técnico de Diseño y Cálculo}
\author{[Nombre del Autor]}
\date{\today}

\begin{document}

% ======================================
% PORTADA / TITLE PAGE
% ======================================
\begin{titlepage}
    \centering
    \vspace*{2cm}
    
    {\Large \textbf{[NOMBRE DE LA INSTITUCIÓN]} \par}
    \vspace{1cm}
    {\large [Facultad/Departamento] \par}
    \vspace{2cm}
    
    {\huge \textbf{Análisis Térmico y de Transferencia de Calor} \par}
    \vspace{0.5cm}
    {\Large \textbf{Informe Técnico de Diseño y Cálculo} \par}
    \vspace{2cm}
    
    {\large \textbf{Por:} \par}
    {\large [Nombre del Autor] \par}
    \vspace{1cm}
    
    {\large \textbf{Curso:} [Nombre del Curso] \par}
    {\large \textbf{Profesor:} [Nombre del Profesor] \par}
    \vspace{2cm}
    
    {\large \today \par}
    
    \vfill
\end{titlepage}

% ======================================
% DESCRIPCIÓN DEL PROBLEMA
% ======================================
\newpage
\section{Descripción del Problema}

[Aquí se debe describir detalladamente el problema de transferencia de calor a analizar. Incluir las condiciones iniciales, el sistema a estudiar, y los objetivos del análisis.]

\vspace{1cm}
% Espacio reservado para diagramas del problema
\begin{figure}[H]
    \centering
    % \includegraphics[width=0.8\textwidth]{diagrama_problema.png}
    \fbox{\parbox{0.8\textwidth}{\centering [ESPACIO RESERVADO PARA DIAGRAMA DEL PROBLEMA] \\ \vspace{3cm}}}
    \caption{Diagrama esquemático del sistema a analizar}
    \label{fig:diagrama_problema}
\end{figure}

\subsection{Objetivos}
\begin{itemize}
    \item [Objetivo específico 1]
    \item [Objetivo específico 2]
    \item [Objetivo específico 3]
\end{itemize}

% ======================================
% TABLA DE DATOS
% ======================================
\newpage
\section{Tabla de Datos}

A continuación se presentan los datos necesarios para el análisis térmico:

\begin{table}[H]
    \centering
    \caption{Datos de entrada para el análisis térmico}
    \label{tab:datos_entrada}
    \begin{tabular}{lcc}
        \toprule
        \textbf{Parámetro} & \textbf{Valor} & \textbf{Unidades} \\
        \midrule
        [Parámetro 1] & [Valor] & [Unidad] \\
        [Parámetro 2] & [Valor] & [Unidad] \\
        [Parámetro 3] & [Valor] & [Unidad] \\
        [Parámetro 4] & [Valor] & [Unidad] \\
        [Parámetro 5] & [Valor] & [Unidad] \\
        \bottomrule
    \end{tabular}
\end{table}

\vspace{1cm}
% Espacio reservado para gráficos de datos
\begin{figure}[H]
    \centering
    % \includegraphics[width=0.8\textwidth]{grafico_datos.png}
    \fbox{\parbox{0.8\textwidth}{\centering [ESPACIO RESERVADO PARA GRÁFICOS DE DATOS] \\ \vspace{3cm}}}
    \caption{Representación gráfica de los datos de entrada}
    \label{fig:datos_graficos}
\end{figure}

% ======================================
% SUPUESTOS
% ======================================
\newpage
\section{Supuestos}

Para el desarrollo del análisis térmico se consideran los siguientes supuestos:

\begin{enumerate}
    \item [Supuesto 1: Describir el primer supuesto]
    \item [Supuesto 2: Describir el segundo supuesto]
    \item [Supuesto 3: Describir el tercer supuesto]
    \item [Supuesto 4: Describir el cuarto supuesto]
    \item [Supuesto 5: Describir el quinto supuesto]
\end{enumerate}

\subsection{Justificación de los Supuestos}
[Aquí se debe justificar cada uno de los supuestos realizados, explicando por qué son válidos para el problema específico y qué impacto tienen en los resultados.]

% ======================================
% MODELO Y ECUACIONES
% ======================================
\newpage
\section{Modelo y Ecuaciones}

\subsection{Modelo Físico}
[Descripción del modelo físico utilizado para representar el problema de transferencia de calor.]

% Espacio reservado para diagrama del modelo
\begin{figure}[H]
    \centering
    % \includegraphics[width=0.8\textwidth]{modelo_fisico.png}
    \fbox{\parbox{0.8\textwidth}{\centering [ESPACIO RESERVADO PARA DIAGRAMA DEL MODELO FÍSICO] \\ \vspace{3cm}}}
    \caption{Modelo físico del sistema}
    \label{fig:modelo_fisico}
\end{figure}

\subsection{Ecuaciones Fundamentales}

\subsubsection{Ecuación de Conservación de Energía}
\begin{equation}
    \dot{E}_{in} - \dot{E}_{out} + \dot{E}_{gen} = \dot{E}_{st}
    \label{eq:conservacion_energia}
\end{equation}

\subsubsection{Ecuación de Conducción de Calor}
\begin{equation}
    \dot{q} = -kA\frac{dT}{dx}
    \label{eq:conduccion}
\end{equation}

\subsubsection{Ecuación de Convección}
\begin{equation}
    \dot{q} = hA(T_s - T_{\infty})
    \label{eq:conveccion}
\end{equation}

\subsubsection{Ecuación de Radiación}
\begin{equation}
    \dot{q} = \epsilon\sigma A(T_s^4 - T_{sur}^4)
    \label{eq:radiacion}
\end{equation}

\subsection{Ecuaciones Específicas del Problema}
[Aquí se deben desarrollar las ecuaciones específicas derivadas para el problema particular, mostrando el desarrollo matemático paso a paso.]

% Espacio reservado para diagramas de ecuaciones
\begin{figure}[H]
    \centering
    % \includegraphics[width=0.8\textwidth]{diagrama_ecuaciones.png}
    \fbox{\parbox{0.8\textwidth}{\centering [ESPACIO RESERVADO PARA DIAGRAMAS DE ECUACIONES] \\ \vspace{3cm}}}
    \caption{Representación esquemática de las ecuaciones del modelo}
    \label{fig:esquema_ecuaciones}
\end{figure}

% ======================================
% METODOLOGÍA EN EES
% ======================================
\newpage
\section{Metodología en EES (Engineering Equation Solver)}

\subsection{Configuración del Problema en EES}
[Descripción de cómo se configuró el problema en EES, incluyendo las variables definidas y las ecuaciones implementadas.]

\subsection{Código EES}
\begin{verbatim}
// Código EES para el análisis térmico
// [Aquí se debe incluir el código completo de EES]

// Definición de variables
T_1 = [valor]    // Temperatura inicial [K]
T_2 = [valor]    // Temperatura final [K]
k = [valor]      // Conductividad térmica [W/m·K]
h = [valor]      // Coeficiente de convección [W/m²·K]

// Ecuaciones del modelo
// [Incluir las ecuaciones implementadas en EES]

\end{verbatim}

\subsection{Validación del Modelo}
[Descripción de cómo se validó el modelo en EES, incluyendo verificaciones de balance de energía y comparaciones con soluciones analíticas cuando sea posible.]

% ======================================
% RESULTADOS BASE
% ======================================
\newpage
\section{Resultados Base}

\subsection{Resultados Principales}
[Presentación de los resultados principales obtenidos del análisis térmico.]

\begin{table}[H]
    \centering
    \caption{Resultados principales del análisis térmico}
    \label{tab:resultados_principales}
    \begin{tabular}{lcc}
        \toprule
        \textbf{Variable} & \textbf{Valor} & \textbf{Unidades} \\
        \midrule
        [Variable 1] & [Valor] & [Unidad] \\
        [Variable 2] & [Valor] & [Unidad] \\
        [Variable 3] & [Valor] & [Unidad] \\
        [Variable 4] & [Valor] & [Unidad] \\
        \bottomrule
    \end{tabular}
\end{table}

% Espacio reservado para gráficos de resultados
\begin{figure}[H]
    \centering
    % \includegraphics[width=0.8\textwidth]{resultados_base.png}
    \fbox{\parbox{0.8\textwidth}{\centering [ESPACIO RESERVADO PARA GRÁFICOS DE RESULTADOS BASE] \\ \vspace{3cm}}}
    \caption{Distribución de temperatura en el sistema}
    \label{fig:resultados_temperatura}
\end{figure}

\subsection{Análisis de los Resultados}
[Interpretación y análisis de los resultados obtenidos, explicando su significado físico y relevancia para el problema.]

% ======================================
% ANÁLISIS PARAMÉTRICO
% ======================================
\newpage
\section{Análisis Paramétrico}

\subsection{Variables de Estudio}
[Identificación de las variables clave para el análisis paramétrico y justificación de su selección.]

\subsection{Rangos de Variación}
\begin{table}[H]
    \centering
    \caption{Rangos de variación para el análisis paramétrico}
    \label{tab:rangos_parametrico}
    \begin{tabular}{lccc}
        \toprule
        \textbf{Parámetro} & \textbf{Valor Mínimo} & \textbf{Valor Máximo} & \textbf{Incremento} \\
        \midrule
        [Parámetro 1] & [Min] & [Max] & [Inc] \\
        [Parámetro 2] & [Min] & [Max] & [Inc] \\
        [Parámetro 3] & [Min] & [Max] & [Inc] \\
        \bottomrule
    \end{tabular}
\end{table}

% Espacio reservado para gráficos del análisis paramétrico
\begin{figure}[H]
    \centering
    % \includegraphics[width=0.8\textwidth]{analisis_parametrico.png}
    \fbox{\parbox{0.8\textwidth}{\centering [ESPACIO RESERVADO PARA GRÁFICOS DEL ANÁLISIS PARAMÉTRICO] \\ \vspace{3cm}}}
    \caption{Variación de la temperatura con respecto a los parámetros clave}
    \label{fig:analisis_parametrico}
\end{figure}

\subsection{Sensibilidad de Variables}
[Análisis de la sensibilidad del sistema a las variaciones de los parámetros estudiados.]

% ======================================
% GRÁFICOS
% ======================================
\newpage
\section{Gráficos}

Esta sección presenta una compilación completa de todos los gráficos generados durante el análisis.

\subsection{Gráficos de Distribución de Temperatura}
% Espacio reservado para múltiples gráficos
\begin{figure}[H]
    \centering
    % \includegraphics[width=0.8\textwidth]{grafico_temp_1.png}
    \fbox{\parbox{0.8\textwidth}{\centering [GRÁFICO 1: DISTRIBUCIÓN DE TEMPERATURA] \\ \vspace{2.5cm}}}
    \caption{Distribución de temperatura - Caso 1}
    \label{fig:temp_caso1}
\end{figure}

\begin{figure}[H]
    \centering
    % \includegraphics[width=0.8\textwidth]{grafico_temp_2.png}
    \fbox{\parbox{0.8\textwidth}{\centering [GRÁFICO 2: DISTRIBUCIÓN DE TEMPERATURA] \\ \vspace{2.5cm}}}
    \caption{Distribución de temperatura - Caso 2}
    \label{fig:temp_caso2}
\end{figure}

\subsection{Gráficos de Flujo de Calor}
\begin{figure}[H]
    \centering
    % \includegraphics[width=0.8\textwidth]{grafico_flujo.png}
    \fbox{\parbox{0.8\textwidth}{\centering [GRÁFICO 3: FLUJO DE CALOR] \\ \vspace{2.5cm}}}
    \caption{Variación del flujo de calor}
    \label{fig:flujo_calor}
\end{figure}

\subsection{Gráficos Comparativos}
\begin{figure}[H]
    \centering
    % \includegraphics[width=0.8\textwidth]{grafico_comparativo.png}
    \fbox{\parbox{0.8\textwidth}{\centering [GRÁFICO 4: COMPARACIÓN DE RESULTADOS] \\ \vspace{2.5cm}}}
    \caption{Comparación entre diferentes configuraciones}
    \label{fig:comparativo}
\end{figure}

% ======================================
% COMPARACIÓN LINEALIZADA
% ======================================
\newpage
\section{Comparación Linealizada}

\subsection{Metodología de Linealización}
[Descripción del método utilizado para linealizar el problema no lineal y las aproximaciones realizadas.]

\subsection{Comparación de Resultados}
\begin{table}[H]
    \centering
    \caption{Comparación entre resultados no lineales y linealizados}
    \label{tab:comparacion_lineal}
    \begin{tabular}{lccc}
        \toprule
        \textbf{Variable} & \textbf{Modelo No Lineal} & \textbf{Modelo Linealizado} & \textbf{Error (\%)} \\
        \midrule
        [Variable 1] & [Valor 1] & [Valor 2] & [Error] \\
        [Variable 2] & [Valor 1] & [Valor 2] & [Error] \\
        [Variable 3] & [Valor 1] & [Valor 2] & [Error] \\
        \bottomrule
    \end{tabular}
\end{table}

\subsection{Análisis de Validez}
[Análisis de la validez de la aproximación lineal y determinación de los rangos donde es aplicable.]

% Espacio reservado para gráficos de comparación
\begin{figure}[H]
    \centering
    % \includegraphics[width=0.8\textwidth]{comparacion_lineal.png}
    \fbox{\parbox{0.8\textwidth}{\centering [GRÁFICO: COMPARACIÓN LINEAL VS NO LINEAL] \\ \vspace{3cm}}}
    \caption{Comparación entre modelos lineal y no lineal}
    \label{fig:comparacion_lineal}
\end{figure}

% ======================================
% RECOMENDACIONES DE DISEÑO
% ======================================
\newpage
\section{Recomendaciones de Diseño}

\subsection{Parámetros Óptimos}
Basándose en los resultados del análisis térmico y paramétrico, se recomiendan los siguientes parámetros de diseño:

\begin{enumerate}
    \item \textbf{[Parámetro 1]:} [Valor recomendado] - [Justificación]
    \item \textbf{[Parámetro 2]:} [Valor recomendado] - [Justificación]
    \item \textbf{[Parámetro 3]:} [Valor recomendado] - [Justificación]
\end{enumerate}

\subsection{Consideraciones de Seguridad}
[Factores de seguridad recomendados y consideraciones especiales para la implementación del diseño.]

\subsection{Limitaciones del Análisis}
[Identificación de las limitaciones del análisis realizado y áreas donde se requiere mayor investigación.]

\subsection{Trabajos Futuros}
[Sugerencias para estudios adicionales que puedan mejorar el diseño o ampliar el análisis.]

% ======================================
% REFERENCIAS
% ======================================
\newpage
\bibliographystyle{apacite}
\bibliography{referencias}

\end{document}